\chapter{BigQuery ML: Machine Learning in SQL}
\index{BigQuery ML}\index{BQML}\index{ARIMA\_PLUS}\index{logistic regression}\index{ML.EVALUATE}\index{churn prediction}

\begin{objectives}
\begin{itemize}
    \item Train regression, classification, and clustering models in pure SQL where the data already lives.
    \item Forecast time series with \texttt{ARIMA\_PLUS} and detect anomalies from residuals.
    \item Evaluate models with \texttt{ML.EVALUATE}, read ROC and precision/recall, and detect overfitting.
    \item Tune hyperparameters and understand BQML's automatic preprocessing and encoding.
    \item Train a logistic-regression churn model on Google Analytics sample data and evaluate its ROC curve.
    \item Architect a pure-SQL forecasting system for daily inventory needs across one thousand stores.
\end{itemize}
\end{objectives}

\begin{crosscloud}
In-warehouse machine learning eliminates the extract-train-reload loop; every major warehouse now ships a variant.

\vspace{2mm}
{\footnotesize
\begin{tabular}{@{}p{2.8cm}p{3.0cm}p{3.0cm}p{3.0cm}@{}} \\
\toprule
\textbf{Capability} & \textbf{GCP BigQuery ML} & \textbf{Snowflake} & \textbf{Azure / Fabric} \\
\midrule
SQL-native training & CREATE MODEL & Snowpark ML / ML functions & Synapse/Fabric ML (Spark MLlib) \\
Forecasting & ARIMA\_PLUS & Forecasting ML function & AutoML forecasting \\
In-warehouse inference & ML.PREDICT & Snowpark ML inference & Synapse predictions \\
Python escape hatch & Remote functions / Vertex AI & Snowpark Python & Spark notebooks \\
Best fit & SQL-fluent teams, tabular data & SQL/Python teams & Fabric-native teams \\
\bottomrule
\end{tabular}}

\vspace{1mm}
\textbf{MongoDB Atlas} users export to analytical stores for ML; \textbf{Redshift ML} wraps SageMaker behind SQL similarly to BQML's managed training. The pattern's shared virtue: governance, lineage, and access control of the warehouse extend automatically to training data and models.
\end{crosscloud}

\begin{keyterms}
\textbf{BQML}: BigQuery ML---model training, evaluation, and prediction expressed in GoogleSQL against warehouse tables.
\textbf{CREATE MODEL}: the DDL statement that trains a model over a query result.
\textbf{ARIMA\_PLUS}: BQML's managed time-series forecaster handling seasonality, holidays, and anomalies automatically.
\textbf{ML.EVALUATE}: the function returning model quality metrics on held-out or supplied data.
\textbf{ROC AUC}: area under the receiver operating characteristic curve; a threshold-independent classifier quality measure.
\textbf{Feature preprocessing}: BQML's automatic one-hot encoding of categoricals and standardization of numerics.
\end{keyterms}

\section{Week 17 Roadmap}
Week~17 opens Part~V with the fastest path from data to model: no exports, no notebooks, no infrastructure. Monday surveys BQML's model catalog. Tuesday forecasts with \texttt{ARIMA\_PLUS}. Wednesday evaluates and tunes. Thursday trains a churn classifier on the Google Analytics sample dataset. Friday architects thousand-store inventory forecasting in pure SQL.

\begin{definitionbox}
\textbf{BigQuery ML} executes machine learning inside the warehouse: \texttt{CREATE MODEL} trains over a SQL query, \texttt{ML.EVALUATE} scores it, \texttt{ML.PREDICT} serves it---with preprocessing, train/test splitting, and hyperparameter defaults managed by the service \cite{googlecloudbqmldocs}.
\end{definitionbox}

\section{Monday: The BQML Catalog---Regression, Classification, Clustering}
\index{BQML!model types}\index{CREATE MODEL}

\subsection{Model Types}
The practical catalog: \textbf{linear and logistic regression} for numeric and binary targets; \textbf{boosted trees} (XGBoost-style) and \textbf{deep neural networks} for nonlinear tabular problems; \textbf{k-means} for segmentation; \textbf{matrix factorization} for recommendations; \textbf{ARIMA\_PLUS} for time series; and \textbf{AutoML Tables} integration when you want BQML syntax over Vertex's AutoML engine. Choosing is mostly standard ML judgment: start with the interpretable baseline (logistic regression), graduate to boosted trees when interactions matter, reach for DNNs last in SQL.

\begin{lstlisting}[language=SQL,caption={Training a churn classifier in pure SQL.}]
CREATE OR REPLACE MODEL analytics.churn_model
OPTIONS(model_type = 'LOGISTIC_REG',
        input_label_cols = ['will_churn'],
        auto_class_weights = TRUE)
AS
SELECT
  totals.visits,
  totals.pageviews,
  totals.timeonsite,
  geoNetwork.country,
  device.deviceCategory,
  IF(totals.transactions IS NULL, 0, 1) AS will_churn
FROM `bigquery-public-data.google_analytics_sample.ga_sessions_*`
WHERE _TABLE_SUFFIX BETWEEN '20160801' AND '20170630';
\end{lstlisting}

\subsection{Why Train Where Data Lives}
Moving data to a training silo costs pipeline code, freshness, and governance: the extract drifts from the warehouse, and the warehouse's access controls no longer apply. In-warehouse training keeps one copy, one permission model, one lineage graph---and makes retraining a scheduled SQL query (Chapter~24 makes it a Dataform model). The limits are equally real: exotic architectures, GPU-heavy training, and non-tabular data belong to Vertex AI (Chapters 18--19).

\begin{pitfall}
\texttt{CREATE MODEL} reads the whole training query every run. Point it at a four-year unfiltered scan and you pay for it repeatedly. Materialize a curated training table (partitioned, clustered) first, train from that, and version it so models are reproducible.
\end{pitfall}

\section{Tuesday: Time-Series Forecasting with ARIMA\_PLUS}
\index{time series}\index{ARIMA\_PLUS!options}\index{anomaly detection}

\subsection{Managed Forecasting}
\texttt{ARIMA\_PLUS} automates the classical ARIMA workflow---differencing, seasonal decomposition, holiday effects, hyperparameter search---behind \texttt{OPTIONS(model\_type='ARIMA\_PLUS', time\_series\_timestamp\_col=..., time\_series\_data\_col=...)}. With \texttt{time\_series\_id\_col} it trains one model per series (per store, per SKU) in a single statement: the thousand-store forecast is one query, not one thousand jobs.

\begin{lstlisting}[language=SQL,caption={Per-store daily demand forecasting and forecast inspection.}]
CREATE OR REPLACE MODEL analytics.store_demand
OPTIONS(model_type = 'ARIMA_PLUS',
        time_series_timestamp_col = 'sale_date',
        time_series_data_col = 'units_sold',
        time_series_id_col = 'store_id',
        holiday_region = 'US')
AS SELECT sale_date, store_id, units_sold
FROM analytics.daily_store_sales;

SELECT store_id, forecast_timestamp, forecast_value,
       prediction_interval_lower_bound, prediction_interval_upper_bound
FROM ML.FORECAST(MODEL analytics.store_demand,
                 STRUCT(30 AS horizon, 0.9 AS confidence_level))
ORDER BY store_id, forecast_timestamp;
\end{lstlisting}

\subsection{Anomaly Detection from Residuals}
\texttt{ML.DETECT\_ANOMALIES} flags points whose observed values fall outside the model's prediction interval---for ARIMA\_PLUS, anomalies in the \emph{history} (a demand spike the seasonality cannot explain) and, applied to new data, anomalies in \emph{operations} (today's sales diverging from forecast). The same primitive powers data-quality monitoring: forecast the row-count time series of a critical table and alert when reality leaves the interval.

\begin{tipbox}
Feed ARIMA\_PLUS clean, gap-filled daily series: missing dates become implicit zeros or gaps depending on options, and both distort seasonality. A \texttt{GENERATE\_DATE\_ARRAY} spine left-joined to actuals is the standard preprocessing idiom.
\end{tipbox}

\section{Wednesday: Evaluation and Hyperparameter Tuning}
\index{ML.EVALUATE!metrics}\index{ROC curve}\index{precision and recall}\index{overfitting}

\subsection{Reading ML.EVALUATE}
For classifiers, \texttt{ML.EVALUATE} returns precision, recall, accuracy, F1, log loss, and ROC AUC. \texttt{ML.ROC\_CURVE} returns per-threshold true/false-positive rates for threshold selection. For imbalanced targets---churn, fraud---accuracy is a liar (predicting ``no churn'' everywhere scores high); precision/recall at the operating threshold and ROC AUC are the honest metrics. Evaluate on a time-split holdout for temporal problems: random splits leak future behavior into training and inflate every metric.

\begin{table}[H]
\centering
\small
\begin{tabular}{@{}p{3.4cm}p{4.6cm}p{4.6cm}@{}} \\
\toprule
\textbf{Metric} & \textbf{Question answered} & \textbf{Churn-model caution} \\
\midrule
ROC AUC & Ranking quality across all thresholds & Threshold-free; hides operating-point pain \\
Precision & Of predicted churners, how many churn? & Low precision floods retention campaigns \\
Recall & Of true churners, how many caught? & Low recall misses saveable customers \\
Accuracy & Overall correctness & Misleading under class imbalance \\
Log loss & Probability calibration quality & Penalizes confident wrongness; tuning target \\
\bottomrule
\end{tabular}
\caption{Classifier evaluation metrics and their churn-prediction reading.}
\label{tab:ch17-eval-metrics}
\end{table}

\subsection{Tuning and Overfitting}
BQML supports hyperparameter tuning (\texttt{num\_trials}, \texttt{hparam\_tuning\_objectives}) over learning rate, regularization, and architecture knobs, reporting trials in \texttt{ML.TRIAL\_INFO}. Detect overfitting by comparing \texttt{ML.EVALUATE} on training versus held-out data: a wide gap (training AUC 0.99, holdout 0.78) says the model memorized; regularize, simplify features, or gather data. With \texttt{auto\_class\_weights=TRUE}, verify the implied threshold on the ROC curve rather than trusting the default 0.5 cutoff.

\begin{pitfall}
A model evaluated on the same period it trained on is a horoscope, not a forecast. For temporal problems, always hold out the most recent period and evaluate there---and repeat the evaluation after every data-pipeline change upstream, because feature drift silently invalidates yesterday's metrics.
\end{pitfall}

\section{Thursday Hands-On Project: Churn Prediction on Google Analytics Data}
\index{hands-on project}\index{Google Analytics sample}\index{logistic regression!churn}
This lab trains a logistic-regression churn model on the public Google Analytics sample dataset and evaluates its ROC curve, entirely in SQL.

\subsection{Train}
Run Wednesday's \texttt{CREATE MODEL} against \texttt{bigquery-public-data.google\_analytics\_sample.ga\_sessions\_*}, using the first eleven months for training. The label flags visitors with transactions; treat the exercise as churn-propensity scoring (high-transaction propensity as the positive class for a clean binary).

\subsection{Evaluate and Predict}
\begin{lstlisting}[language=SQL,caption={Evaluating the churn model and reading its ROC curve.}]
SELECT *
FROM ML.EVALUATE(MODEL analytics.churn_model, (
  SELECT totals.visits, totals.pageviews, totals.timeonsite,
         geoNetwork.country, device.deviceCategory,
         IF(totals.transactions IS NULL, 0, 1) AS will_churn
  FROM `bigquery-public-data.google_analytics_sample.ga_sessions_*`
  WHERE _TABLE_SUFFIX BETWEEN '20170701' AND '20170801'));

SELECT threshold, recall, true_positives, false_positives, true_negatives
FROM ML.ROC_CURVE(MODEL analytics.churn_model, (
  SELECT totals.visits, totals.pageviews, totals.timeonsite,
         geoNetwork.country, device.deviceCategory,
         IF(totals.transactions IS NULL, 0, 1) AS will_churn
  FROM `bigquery-public-data.google_analytics_sample.ga_sessions_*`
  WHERE _TABLE_SUFFIX BETWEEN '20170701' AND '20170801'))
ORDER BY threshold;

SELECT fullVisitorId, predicted_will_churn_probs
FROM ML.PREDICT(MODEL analytics.churn_model, (
  SELECT fullVisitorId, totals.visits, totals.pageviews, totals.timeonsite,
         geoNetwork.country, device.deviceCategory
  FROM `bigquery-public-data.google_analytics_sample.ga_sessions_20170801`))
ORDER BY (SELECT prob FROM UNNEST(predicted_will_churn_probs) WHERE label = 1) DESC
LIMIT 100;
\end{lstlisting}

\subsection*{Deliverables}
\begin{itemize}
    \item The training, evaluation, and prediction SQL committed with a README on the time split.
    \item The \texttt{ML.EVALUATE} output and the chosen operating threshold with justification.
    \item A short paragraph comparing boosted-tree and logistic-regression results on the same data.
\end{itemize}

\subsection*{Acceptance Criteria}
\begin{itemize}
    \item Training and evaluation use disjoint time periods.
    \item Metrics reported include ROC AUC and precision/recall at a stated threshold, not accuracy alone.
    \item The prediction query produces calibrated probability output, not just labels.
\end{itemize}

\subsection*{Stretch Goals}
\begin{itemize}
    \item Retrain with \texttt{BOOSTED\_TREE\_CLASSIFIER} and \texttt{ML.EXPLAIN\_PREDICT} for feature attributions.
    \item Run a ten-trial hyperparameter search and report the trial table.
    \item Schedule weekly retraining via a scheduled query and log metrics per run.
\end{itemize}

\section{Friday Use Case and Architecture: SQL-Only Forecasting for 1,000 Stores}
\index{inventory forecasting}\index{SQL pipelines}\index{MLOps!BQML}

\subsection{Scenario}
A grocery chain must forecast daily unit demand per store per category---one thousand stores, fifty categories---for replenishment planning. The analytics team is SQL-fluent; there is no ML platform team; the forecast must refresh nightly and feed the ordering system by 05:00.

\subsection{Architecture}
The whole system is scheduled SQL: a Dataform or scheduled-query pipeline (1) rebuilds the gap-filled daily series per store-category from the sales replica, (2) retrains \texttt{ARIMA\_PLUS} weekly with \texttt{time\_series\_id\_col} covering store-category pairs, (3) nightly runs \texttt{ML.FORECAST} with a seven-day horizon into a forecast table partitioned by forecast date, (4) applies business rules (min/max stock, case-pack rounding) in SQL, and (5) exports the order file to the ERP landing bucket. Monitoring is itself a forecast: \texttt{ML.DETECT\_ANOMALIES} over realized-versus-forecast error, paged when sustained.

\begin{figure}[H]
    \centering
    \resizebox{0.95\textwidth}{!}{%
    \begin{tikzpicture}[
        db/.style={cylinder, shape border rotate=90, aspect=0.25, draw=primary, thick, fill=primary!10, text width=2.9cm, minimum height=1.2cm, align=center, font=\small\bfseries},
        svc/.style={rectangle, rounded corners, draw=primary, thick, fill=primary!6, text width=3.0cm, minimum height=1.0cm, align=center, font=\small\bfseries},
        outbox/.style={rectangle, rounded corners, draw=codegreen!60!black, thick, fill=green!7, text width=2.9cm, minimum height=1.0cm, align=center, font=\small\bfseries},
        arrow/.style={-{Stealth[length=2.5mm]}, draw=primary, thick}
    ]
        \node[db] (sales) at (0,0) {Sales replica\\(CDC, Ch. 11)};
        \node[svc] (series) at (4.2,0) {Scheduled SQL\\series build\\(gap-filled)};
        \node[svc] (model) at (8.4,0) {ARIMA\_PLUS\\weekly retrain\\50k series};
        \node[svc] (forecast) at (12.4,0) {Nightly\\ML.FORECAST\\+ business rules};
        \node[outbox] (erp) at (15.6,0) {ERP order\\file 05:00};

        \draw[arrow] (sales) -- (series);
        \draw[arrow] (series) -- (model);
        \draw[arrow] (model) -- (forecast);
        \draw[arrow] (forecast) -- (erp);
    \end{tikzpicture}%
    }
    \caption{SQL-only inventory forecasting: CDC replica to nightly per-series forecasts to the ERP.}
    \label{fig:ch17-inventory-forecast}
\end{figure}

\begin{tipbox}
Fifty thousand series in one \texttt{CREATE MODEL} is normal for ARIMA\_PLUS, but watch the long tail: stores with sparse histories produce wide intervals. Segment by data sufficiency and route thin series to a simpler fallback (moving average) rather than trusting a heroic model.
\end{tipbox}

\begin{pitfall}
A forecast table without a forecast-generation timestamp is unusable for accountability: you can never reconstruct ``what did we predict on Tuesday for Friday.'' Always persist \texttt{forecast\_for\_date}, \texttt{generated\_at}, and model version with every batch.
\end{pitfall}

\section{Interview Q\&A}
\subsection{Concepts and Trade-offs}
\begin{enumerate}
    \item \textbf{Q: Which BQML model type fits customer churn prediction?}\\
    A: Logistic regression as the interpretable baseline; boosted trees when nonlinear feature interactions justify the opacity.

    \item \textbf{Q: What does ARIMA\_PLUS automate compared to classic ARIMA?}\\
    A: Differencing, seasonality, holiday effects, hyperparameter search, and per-series parallelism via \texttt{time\_series\_id\_col}.

    \item \textbf{Q: Which ML.EVALUATE metrics matter most for imbalanced churn?}\\
    A: ROC AUC for ranking quality plus precision and recall at the chosen operating threshold; accuracy misleads under imbalance.

    \item \textbf{Q: What is the advantage of training where the data already lives?}\\
    A: No extract pipelines, one permission model, full lineage, and retraining as a scheduled SQL query.

    \item \textbf{Q: How do you detect overfitting with BQML evaluation output?}\\
    A: Compare ML.EVALUATE on training versus time-held-out data; a large quality gap indicates memorization.
\end{enumerate}

\section{Further Reading}
Read the BigQuery ML introduction and model reference for model types, options, and evaluation functions \cite{googlecloudbqmldocs}. The BigQuery documentation covers scheduled queries that productionize retraining \cite{googlecloudbigquerydocs}. Vertex AI documentation previews the custom-training path for when BQML's catalog is insufficient \cite{googlecloudvertexaidocs}, and Chapter~24's Dataform reference covers SQL pipeline discipline \cite{googleclouddataformdocs}.

\section*{Key Takeaways}
\begin{itemize}
    \item BQML moves ML to the data: governance, freshness, and lineage come free.
    \item \texttt{ARIMA\_PLUS} turns thousand-series forecasting into one query---given clean, gap-filled inputs.
    \item Evaluate on time-held-out data; report threshold-aware metrics for imbalanced problems.
    \item Materialize curated training tables; never train over raw multi-year scans.
    \item Persist forecast provenance (generated-at, version) with every batch.
    \item Know the boundary: exotic architectures and non-tabular data belong to Vertex AI.
\end{itemize}

\section{Exercises}
\begin{enumerate}
    \item \textbf{(Conceptual)} When does BQML's in-warehouse training lose to a Vertex AI custom job? Give three concrete cases.
    \item \textbf{(Conceptual)} Explain why random train/test splits inflate metrics for temporal targets.
    \item \textbf{(Hands-on)} Reproduce the Thursday lab, then build the ROC-based threshold analysis for a retention-budget constraint (campaign capacity of the top 5\%).
    \item \textbf{(Hands-on)} Forecast one public time series with ARIMA\_PLUS and verify interval coverage on a held-out month.
    \item \textbf{(Design)} Specify the anomaly-detection layer for the store forecasts: error metrics, alert thresholds, and the human triage loop.
    \item \textbf{(Architecture)} Add a feedback table capturing realized orders versus forecasts, and design the weekly accuracy review query.
\end{enumerate}
